% ID: SEL-D0031-OB-E02
% ESTADO: cotejada. Segunda pasada visual completa y borrador compilado revisados.
% Fuente: Selectividad/matematicas_aplicadas/titular_junio_6 EXAMEN MAT APLIC C S.pdf
% SHA-256: bebc339cfd8f69c7c4fae160fd165eaef5f6bc93cfd6d15cb9ce289b7029d167
% Página: 2; opción B; ejercicio original 2.

La cantidad, $C$, que una entidad bancaria dedica a créditos depende de su liquidez, $x$, según la función
\[
C(x)=\begin{cases}
\dfrac{150+5x}{100} & \text{si }10\leq x\leq50\\[4pt]
\dfrac{200+10x}{25+3x} & \text{si }x>50
\end{cases}
\]
donde $C$ y $x$ están expresadas en miles de euros.
\begin{enumerate}[label=\alph*)]
\item \textbf{(1 punto)} Justifique que $C$ es una función continua.
\item \textbf{(1 punto)} \textquestiondown{}A partir de qué liquidez decrece la cantidad dedicada a créditos? \textquestiondown{}Cuál es el valor máximo de $C$?
\item \textbf{(0.5 puntos)} Calcule la asíntota horizontal e interprétela en el contexto del problema.
\end{enumerate}
